% !TEX program = xelatex
\documentclass[10pt,a4paper]{article}
\usepackage[a4paper,margin=0pt]{geometry}
\usepackage{fontspec}
\usepackage{polyglossia}
\usepackage[dvipsnames]{xcolor}
\usepackage{tikz}
\usepackage{eso-pic}
\usepackage{enumitem}
\usepackage{tabularx}
\usepackage{array}
\usepackage{calc}
\usepackage{graphicx}
\usepackage{textcomp}
\usepackage[hidelinks]{hyperref}
\usepackage[final]{microtype}
\usepackage{ragged2e}

\defaultfontfeatures{Ligatures=TeX}
\setmainfont{Rubik-Regular.ttf}[
  Path = fonts/,
  BoldFont = Rubik-Bold.ttf,
  ItalicFont = Rubik-Italic.ttf,
  BoldItalicFont = Rubik-BoldItalic.ttf
]
\newfontfamily\RubikRegular{Rubik-Regular.ttf}[
  Path = fonts/,
  BoldFont = Rubik-Medium.ttf,
  ItalicFont = Rubik-Italic.ttf,
  BoldItalicFont = Rubik-BoldItalic.ttf
]
\newfontfamily\RubikLight{Rubik-Light.ttf}[
  Path = fonts/,
  ItalicFont = Rubik-Italic.ttf
]
\newfontfamily\TemplateFont{cmunrm.otf}[
  Path = fonts/,
  BoldFont = cmunbx.otf,
  ItalicFont = cmunti.otf,
  BoldItalicFont = cmunbi.otf
]
\newfontfamily\cyrillicfont{Rubik-Regular.ttf}[
  Script = Cyrillic,
  Path = fonts/,
  BoldFont = Rubik-Bold.ttf,
  ItalicFont = Rubik-Italic.ttf,
  BoldItalicFont = Rubik-BoldItalic.ttf
]

\pagestyle{empty}
\setlength{\parindent}{0pt}
\setlength{\parskip}{0pt}
\setlength{\tabcolsep}{0pt}
\setlength{\fboxsep}{0pt}
\setlength{\emergencystretch}{2em}
\raggedbottom
\urlstyle{same}
\hypersetup{colorlinks=true, urlcolor=black, linkcolor=black}

\definecolor{sidebarbg}{HTML}{244869}
\definecolor{accentblue}{HTML}{0B77F3}
\definecolor{textgray}{HTML}{555555}
\definecolor{rulegray}{HTML}{B7B7B7}
\definecolor{sidebarline}{HTML}{DCE6F0}
\definecolor{softwhite}{HTML}{F7FAFC}

\newlength{\sidebarwidth}
\setlength{\sidebarwidth}{0.355\paperwidth}
\newlength{\mainwidth}
\setlength{\mainwidth}{\paperwidth-\sidebarwidth}
\newlength{\sidebarinnerwidth}
\setlength{\sidebarinnerwidth}{\sidebarwidth-2.12cm}
\newlength{\maininnerwidth}
\setlength{\maininnerwidth}{\mainwidth-1.72cm}
\newlength{\pagecontentheight}
\setlength{\pagecontentheight}{\paperheight-1.05cm}
\newlength{\sidebarinnerheight}
\setlength{\sidebarinnerheight}{\pagecontentheight-0.84cm}
\newlength{\maininnerheight}
\setlength{\maininnerheight}{\pagecontentheight-0.74cm}
\newlength{\maindatewidth}
\setlength{\maindatewidth}{2.65cm}
\newlength{\maingap}
\setlength{\maingap}{0.30cm}

\newlength{\avatarheight}
\setlength{\avatarheight}{5.33cm}

\newcommand{\atsbackground}{%
  \AddToShipoutPictureBG*{%
    \AtPageLowerLeft{\color{sidebarbg}\rule{\sidebarwidth}{\paperheight}}%
  }%
}

\newcommand{\TemplateText}[1]{{\TemplateFont #1}}
\newcommand{\MetaLabel}[1]{{\TemplateFont\fontsize{8.55}{8.9}\selectfont #1}}

\newcommand{\SidebarAvatar}{%
  {\centering\includegraphics[height=\avatarheight,keepaspectratio]{data/avatar.png}\par}
  \vspace{0.30cm}%
}
\newcommand{\CVName}[1]{%
  {\centering\TemplateFont\fontsize{25.0}{24.4}\selectfont\bfseries\MakeUppercase{#1}\par}
  \vspace{0.66cm}
}

\newcommand{\SidebarSection}[1]{%
  {\TemplateFont\fontsize{12.2}{12.7}\selectfont\MakeUppercase{#1}\par}
  \vspace{-0.30cm}
  {\color{sidebarline}\rule{\linewidth}{0.48pt}\par}
  \vspace{0.10cm}
}

\newcommand{\SidebarRoleTitle}[1]{{\RubikRegular\fontsize{9.9}{10.55}\selectfont\bfseries #1\par}}
\newcommand{\SidebarOrg}[1]{{\RubikRegular\fontsize{9.15}{9.8}\selectfont #1\par}}
\newcommand{\SidebarDate}[1]{{\RubikRegular\fontsize{8.88}{9.4}\selectfont #1\par}}
\newcommand{\SidebarBody}[1]{{\RubikLight\fontsize{8.95}{9.7}\selectfont #1\par}}

\newenvironment{sidebarbullets}{%
  \begin{itemize}[leftmargin=1.08em,topsep=0.08em,itemsep=0.05em,parsep=0pt,partopsep=0pt]
    \RubikLight\fontsize{8.72}{9.55}\selectfont\color{softwhite}
}{\end{itemize}}

\newcommand{\SidebarEntrySep}{\vspace{0.30cm}}
\newcommand{\SidebarSectionSep}{\vspace{0.68cm}}

\newcommand{\EducationEntry}[4]{%
  {\fontsize{10.15}{10.65}\selectfont\bfseries #1\par}
  \vspace{0.04cm}
  {\fontsize{9.1}{9.65}\selectfont #2\par}
  \vspace{0.05cm}
  {\fontsize{8.9}{9.4}\selectfont #3\ \textbar\ #4\par}
}

\newcommand{\LanguageEntry}[2]{%
  \begin{tabular*}{\linewidth}{@{}>{\fontsize{9.7}{10.15}\selectfont}l@{\extracolsep{\fill}}>{\fontsize{9.1}{9.55}\selectfont}r@{}}
    #1 & #2
  \end{tabular*}\par
}

\newcommand{\SkillLine}[1]{{\fontsize{9.2}{9.75}\selectfont #1\par}}

\newcommand{\InterestTitle}[1]{{\RubikRegular\fontsize{10.35}{10.95}\selectfont\bfseries #1\par\vspace{0.03cm}}}
\newcommand{\InterestBody}[1]{{\RubikLight\fontsize{8.75}{9.6}\selectfont #1\par}}

\newcommand{\TopRole}[1]{{\fontsize{13.8}{14.15}\selectfont\color{accentblue}#1\par}}
\newcommand{\ContactLine}[1]{{\fontsize{8.25}{8.7}\selectfont\color{textgray}#1\par}}
\newcommand{\HeaderGap}{\vspace{0.46cm}}

\newcommand{\MainSection}[1]{%
  {\TemplateFont\fontsize{13.6}{14.1}\selectfont\MakeUppercase{#1}\par}
  \vspace{-0.30cm}
  {\color{rulegray}\rule{\linewidth}{0.46pt}\par}
  \vspace{0.12cm}
}

\newcommand{\SummaryText}[1]{
  {\fontsize{9.9}{10.5}\selectfont\color{black!85}#1\par}
  \vspace{0.35cm}
}

\newcommand{\MainEntry}[5]{%
  \begin{tabular*}{\linewidth}{@{}p{\dimexpr\linewidth-\maindatewidth-\maingap\relax}@{\hspace{\maingap}}>{\raggedleft\arraybackslash}p{\maindatewidth}@{}}
    {\fontsize{10.75}{11.1}\selectfont #1} & {\fontsize{9.2}{9.6}\selectfont #2}
  \end{tabular*}\par
  \vspace{-0.30cm}
  \begin{tabular*}{\linewidth}{@{}p{\dimexpr\linewidth-\maindatewidth-\maingap\relax}@{\hspace{\maingap}}>{\raggedleft\arraybackslash}p{\maindatewidth}@{}}
    {\fontsize{9.9}{10.3}\selectfont\color{accentblue}#3} & {\fontsize{8.95}{9.35}\selectfont\raggedleft\color{textgray}#4}
  \end{tabular*}\par
  \vspace{0.02cm}
  #5
  \vspace{0.35cm}
}

\newenvironment{mainbullets}{%
  \begin{itemize}[leftmargin=1.0em,topsep=0.02em,itemsep=0.04em,parsep=0pt,partopsep=0pt]
    \fontsize{9.3}{9.8}\selectfont\color{black!82}
}{\end{itemize}}

\newcommand{\CoursesSection}[1]{%
  {\TemplateFont\fontsize{13.6}{14.1}\selectfont\MakeUppercase{#1}\par}
  \vspace{-0.30cm}
  {\color{rulegray}\rule{\linewidth}{0.46pt}\par}
  \vspace{0.12cm}
}

\newcommand{\CourseItem}[2]{%
  {\fontsize{9.9}{10.3}\selectfont\color{accentblue}#1\par}
  \vspace{0.02cm}
  {\fontsize{9.15}{9.6}\selectfont\color{black!82}#2\par}
  \vspace{0.30cm}
}

\newcommand{\RightFooter}[1]{\vfill{\hfill\fontsize{8.0}{8.0}\selectfont\color{textgray}#1\par}}
\newcommand{\LeftFooter}[1]{\vfill{\fontsize{8.0}{8.0}\selectfont\color{sidebarline}#1\par}}

\setmainlanguage{russian}
\setotherlanguage{english}

\begin{document}
\atsbackground
\noindent\makebox[\paperwidth][l]{%
\begin{minipage}[t][\pagecontentheight]{\sidebarwidth}
  \vspace*{0.46cm}
  \hspace*{1.06cm}\begin{minipage}[t][\sidebarinnerheight]{\sidebarinnerwidth}
    \color{softwhite}\RaggedRight
    \SidebarAvatar
    \CVName{Битлев\\Роберт}

    \SidebarSection{Достижения}
    \SidebarRoleTitle{Олимпиады по физике\\и информатике}
    \vspace{0.05cm}
    \begin{sidebarbullets}
      \item Победитель и призёр заключительного этапа ВсОШ по физике.
      \item Золотая медаль IJSO в 2023 и 2024 годах.
      \item Победитель олимпиады «Технокубок» в 2025 году.
      \item Победитель и призёр регионального этапа ВсОШ по информатике.
    \end{sidebarbullets}

    \SidebarSectionSep
    \SidebarSection{Образование}
    \EducationEntry{AI360: Математика\\машинного обучения}{Санкт-Петербургский государственный университет}{2025--2029}{Санкт-Петербург}
    \SidebarEntrySep
    \EducationEntry{Среднее общее образование}{Лицей №\,533}{2016--2025}{Средний балл: 5.00}

    \SidebarSectionSep
    \SidebarSection{Языки}
    \LanguageEntry{Английский}{C1}
    \SidebarEntrySep
    \LanguageEntry{Немецкий}{B2--C1}
    \SidebarEntrySep
    \LanguageEntry{Русский}{Родной}

    \SidebarSectionSep
    \SidebarSection{Навыки}
    \SkillLine{Python \textperiodcentered\ SQL \textperiodcentered\ C++}
    \SidebarEntrySep
    \SkillLine{Машинное обучение \textperiodcentered\ Анализ данных}
    \SidebarEntrySep
    \SkillLine{Алгоритмы и структуры данных}
    \SidebarEntrySep
    \SkillLine{Java \textperiodcentered\ HTML/CSS/JavaScript}

    \SidebarSectionSep
    \SidebarSection{Интересы}
    \InterestTitle{Математика, ML\\и разработка}
    \InterestBody{Интересуюсь машинным обучением, анализом данных, алгоритмическими задачами и прикладной разработкой.}
    \LeftFooter{}
  \end{minipage}
\end{minipage}%
\begin{minipage}[t][\pagecontentheight]{\mainwidth}
  \vspace*{0.70cm}
  \hspace*{0.66cm}\begin{minipage}[t][\maininnerheight]{\maininnerwidth}
    \TopRole{Студент AI360 \textbar\ Математика машинного обучения}
    \vspace{0.09cm}
    \ContactLine{{\MetaLabel{Тел.:}}\ \href{tel:+79215747178}{+7 (921) 574-71-78} \textbar\ {\MetaLabel{Email:}}\ \href{mailto:robertbitlev09@gmail.com}{robertbitlev09@gmail.com}}
    \vspace{0.03cm}
    \ContactLine{{\MetaLabel{GitHub:}}\ \href{https://github.com/robrob09/}{robrob09} \textbar\ {\MetaLabel{Codeforces:}}\ \href{https://codeforces.com/profile/Robrob09}{Robrob09} \textbar\ {\MetaLabel{Telegram:}}\ \href{https://t.me/robrob2009}{@robrob2009}}
    \HeaderGap

    \MainSection{Профиль}
    \SummaryText{Студент программы «AI360: Математика машинного обучения» Факультета математики и компьютерных наук СПбГУ. Интересуюсь машинным обучением, анализом данных и прикладной разработкой; работал над учебными и проектными задачами, связанными с обработкой больших массивов данных, статистическим анализом, прогнозированием и практическими программными инструментами.}

    \vspace{0.24cm}
    \MainSection{Проекты}
    \MainEntry
      {Конструктор резюме}
      {Лето 2025}
      {Статическое веб-приложение с экспортом в PDF}
      {}
      {\begin{mainbullets}
        \item Сверстал конструктор резюме по макету Figma на HTML, CSS и JavaScript: 2 макета, редактирование содержимого в интерфейсе и адаптация под разные экраны в рамках 1-страничного приложения.
        \item Добавил сохранение данных в localStorage и браузерный экспорт в PDF, что позволило редактировать, сохранять, сбрасывать и выгружать текущую версию без backend-части.
      \end{mainbullets}}

    \MainEntry
      {Оценка нагрузки на службу поддержки}
      {Весна 2025}
      {Прогнозирование обращений в техподдержку}
      {}
      {\begin{mainbullets}
        \item Обработал более 600 тыс. обращений и данные продуктовой навигации с помощью Python и SQL; проверил пропуски и дубликаты, собрал метрики пользовательского портрета, повторных обращений и связи поддержки с ростом DAU.
        \item Построил SARIMAX-прогноз недельной нагрузки на 12 недель и логистическую модель вероятности обращения; оформил результаты в виде графиков и рекомендаций для планирования нагрузки и распределения ресурсов.
      \end{mainbullets}}

    \MainEntry
      {Влияние наличия кредитной карты на получение кредита}
      {Весна 2024}
      {Хакатон DANO, Университет ИТМО}
      {}
      {\begin{mainbullets}
        \item Выполнил предобработку и SQL-агрегацию клиентских данных, сравнил распределения по четырём кредитным продуктам и визуализировал зависимости с помощью гистограмм, диаграмм размаха и тепловых карт корреляций.
        \item Выявил факторы выбора кредитного продукта, сформулировал выводы для сегментации аудитории и предложил 3 CRM-сценария: upsell, cross-sell и down-sell.
      \end{mainbullets}}

    \MainEntry
      {Brain Train}
      {2023--2024}
      {Интерактивный тренажёр памяти на Arduino Nano}
      {}
      {\begin{mainbullets}
        \item Спроектировал устройство на Arduino Nano с генерацией случайных LED-последовательностей, OLED-индикацией, системой подсказок и кнопочным вводом для тренировки памяти.
        \item Объединил 4 модуля на Arduino C++ - игровую логику, обработку кнопок, вывод на дисплей и LED-паттерны - и настроил тайминги для более плавного сценария взаимодействия.
      \end{mainbullets}}

    \CoursesSection{Дополнительное образование / Курсы}
    \CourseItem{Т-интенсив по JavaScript}{Т-Образование, лето 2025}
    \CourseItem{Хакатон DANO}{Университет ИТМО, весна 2024}
    \CourseItem{Школа по практическому программированию и анализу данных}{НИУ ВШЭ, 2024--2025}
    \CourseItem{Творческая радиоэлектроника}{Академия цифровых технологий, 2023--2024}
    \CourseItem{Олимпиадное программирование для школьников}{Яндекс Образование, 2023--2025}
    \CourseItem{Алгоритмы и структуры данных}{Т-Поколение, 2021--2022}
    \RightFooter{}
  \end{minipage}
\end{minipage}}
\end{document}
